归并排序（Merge Sort）是「分治」思想最经典的应用之一。
把一个大数组切成两半，分别排序，再把两个有序数组合并成一个。

\subsection{核心思路}

分治的三个阶段：

\begin{enumerate}
    \item \textbf{分解（Divide）}：把数组从中间切开，得到两个子数组
    \item \textbf{解决（Conquer）}：递归地对两个子数组调用归并排序
    \item \textbf{合并（Merge）}：把两个有序的子数组合并成一个有序数组
\end{enumerate}

\subsection{时间复杂度}

因为每次都把数组对半切开，递归树的深度是 \(O(\log n)\)，
每一层合并的总工作量是 \(O(n)\)，所以总时间 \(O(n \log n)\)。
代价是需要一块大小为 \(O(n)\) 的辅助数组 \texttt{B[]}。

\subsection{C 代码 — 合并函数}

核心在于 \texttt{Merge}：把 \texttt{A[low..mid]} 和 \texttt{A[mid+1..high]}
这两段已有序的部分合并为 \texttt{A[low..high]}。

\lstinputlisting[style=cstyle, firstline=1, lastline=34,
    caption=src/sort/merge\_sort.c — Merge 函数（含逆序对计数）]
{../algo/src/sort/merge_sort.c}

这段代码还顺便统计了\textbf{逆序对}（inversion）—— 也就是数组中
`\(i < j\) 但 \(A[i] > A[j]\)` 的情况数。关键一行是：

\begin{lstlisting}[style=cstyle]
count += mid - i + 1;   // 当从右半部分取出元素时
\end{lstlisting}

意思是：当 \texttt{B[j]} 比 \texttt{B[i]} 小，说明从 \texttt{B[i]}
到 \texttt{B[mid]} 这些元素全都是比 \texttt{B[j]} 大的，
每一个都和 \texttt{B[j]} 构成一个逆序对。

\subsection{C 代码 — 递归排序函数与主程序}

\lstinputlisting[style=cstyle, firstline=36, lastline=61,
    caption=src/sort/merge\_sort.c — Mergesort 递归函数 + main]
{../algo/src/sort/merge_sort.c}

注意这段代码使用了\textbf{从 1 开始}的数组下标习惯（\texttt{A[1..n]}），
这是《算法导论》里常用的写法，便于计算 \texttt{mid}。
如果换成从 0 开始，只要把 \texttt{low} 的初始值改成 0、
\texttt{high} 改成 \texttt{n-1} 即可。

\subsection{编译与运行}

\begin{lstlisting}[style=bashstyle]
clang -o merge algo/src/sort/merge_sort.c
./merge
1 2 3 4 5 6 7 8
逆序对个数为：10
\end{lstlisting}

对于测试数组 \{0, 5, 2, 4, 6, 1, 3, 7, 8\}（位置 0 被忽略），
逆序对共 10 对：(5,2), (5,1), (5,3), (2,1), (4,1), (4,3), (6,1), (6,3), 等等。
这给了你一个判断排序实现是否正确的额外指标。
